The Centers for Medicare \& Medicaid Services (CMS) publish aggregate information on the use of and payments for durable medical equipment (DME), devices, and supplies in Original Medicare Part B. The geography-and-service public-use file organizes observations by geographic area, Healthcare Common Procedure Coding System (HCPCS) code, and supplier rental classification, with measures of submitted charges, services, and payments \citep{cmsdme2023}. These aggregated data make it possible to examine the relationship between equipment categories, payment definitions, and regional variation without using beneficiary-level records.

A single pooled prediction score may conceal differences in the reliability of predictions across areas. Randomly partitioning rows permits records from the same geography to occur in both training and testing sets; it therefore answers a different question from predicting payments in an area entirely absent from model training. Methodological research on structured data emphasizes choosing validation schemes appropriate to the intended generalization setting \citep{roberts2017}. Nevertheless, geographic holdout need not always produce a larger numerical error than a random split, particularly when their test-set compositions differ.

A second issue concerns the definition of the prediction target. CMS reports both observed and standardized average Medicare payment amounts. Previous work on Medicare payment standardization explains how geographic and policy-based reimbursement differences can complicate comparisons across regions \citep{odonnell2013}. For DME records, prediction may also depend on the distinction between equipment furnished under rental and non-rental classifications. Furthermore, an HCPCS code absent from the training set cannot be assigned a directly estimated code-specific training mean.

This study investigates the reliability of simple, interpretable payment-prediction baselines across U.S. states and Washington, DC. We evaluate: (i) random versus geographically grouped validation; (ii) known versus unseen HCPCS codes; (iii) actual versus standardized payment targets; and (iv) whether jointly modeling HCPCS and rental classification changes prediction error. Rather than proposing a new learning algorithm, the study develops a reproducible empirical assessment of what conventional aggregate accuracy metrics do and do not establish about geographic prediction reliability.